\bookchapter{الفصل 5: اجعل هاتفك يفهمك أكثر — الإعدادات الأساسية}

\chapterintro{اقضِ بضع دقائق في الإعداد، وستتحسن التجربة بشكل كبير. هذه الإعدادات مثل“ضبط النغمة” للهاتف، لجعله أكثر توافقاً مع عاداتك.}

\booksection{أولاً: الصوت والاهتزاز}

عيّن نغمة وإشعاراً يعجبانك. استخدم وضع الصامت (في الاجتماعات) وعدم الإزعاج (عند النوم) بمرونة، سيصبح العالم هادئاً فوراً.

\begin{itemize}
  \item \textbf{النغمة}: اختر واحدة تسمعها ولا تكون محرجة.
  \item \textbf{وضع الاهتزاز}: التنبيه بالاهتزاز في الوضع الصامت.
  \item \textbf{وضع عدم الإزعاج}: يمكن ضبط استثناءات، مثل السماح بمكالمات جهات الاتصال.
  \item \textbf{فصل مستوى الصوت}: يمكن ضبط مستوى صوت الوسائط والنغمة والمنبه بشكل منفصل.
\end{itemize}

\booksection{ثانياً: العرض والسطوع}

\begin{itemize}
  \item \textbf{السطوع التلقائي}: دع الهاتف يضبط السطوع وفقاً للإضاءة المحيطة، لراحة بالك.
  \item \textbf{وضع حماية العين/العرض الليلي}: يرشّح الضوء الأزرق، لمشاهدة الهاتف ليلاً بشكل أكثر راحة.
  \item \textbf{الوضع الداكن}: يوفّر الطاقة (خاصة لشاشات OLED)، ولا يكون مزعجاً ليلاً. تدعم العديد من التطبيقات الوضع الداكن.
  \item \textbf{معدل تحديث الشاشة}: يمكن ضبطه في بعض هواتف Android، وكلما زاد (مثل 120Hz) كان التمرير أكثر سلاسة.
  \item \textbf{قفل الشاشة التلقائي}: اضبطه على 15 أو 30 ثانية، آمن وموفّر للطاقة.
  \item \textbf{حجم الخط}: هل الخط صغير؟ كبّره قليلاً.
\end{itemize}

\booksection{ثالثاً: الخلفيات والسمات وسطح المكتب}

غيّر خلفية جميلة (ثابتة أو متحركة)، سيتغير مزاجك. تملك هواتف Android عادة متاجر سمات، يمكنها تغيير الأيقونات والأسلوب بالكامل.

\textbf{عناصر سطح المكتب}: ضع الطقس والأجندة والتحكم بالموسيقى على سطح المكتب، للوصول المباشر. يدعم كل من iOS وAndroid ذلك، اضغط مطولاً على مساحة فارغة في سطح المكتب للإضافة.

\booksection{رابعاً: إعدادات الكفاءة: اجعل هاتفك“دماغاً مساعداً” لك}

هذا الجزء يتضمن وظائف“بمجرد إعدادها لا يمكن العودة”، كل منها يرفع كفاءتك اليومية بشكل ملحوظ.

\booksubsection{وضع التركيز/عدم الإزعاج: استعد انتباهك الذي أخذته الإشعارات}

\begin{itemize}
  \item \textbf{الفكرة الأساسية}: ليست"مشاهدة الهاتف أقل"، بل"التركيز وقت التركيز، والتواصل وقت التواصل".
  \item \textbf{توصيات الاستخدام}:\newline · وضع العمل: السماح فقط بجهات الاتصال"المميزة" والتطبيقات الضرورية؛\newline · وضع النوم: تفعيل تلقائي، السماح بالمنبه فقط؛\newline · اللياقة/القراءة/الاجتماعات: التخصيص وفقاً للنشاط الحالي.
  \item \textbf{استخدام متقدم}: استخدم الأتمتة + الموقع/الوقت للتبديل التلقائي بين الأوضاع. مثلاً: يمكن لإجراءات iOS المختصرة ضبط"الدخول التلقائي إلى وضع العمل عند الوصول إلى الشركة".
\end{itemize}

\booksubsection{إدارة الإشعارات: ليتوقف“الاتصال المتتالي” للنقطة الحمراء}

\begin{itemize}
  \item \textbf{إغلاق إشعارات التسويق}: 90\% من الإشعارات لا تحتاجها فعلاً.
  \item \textbf{الترتيب حسب الأهمية}: ضع التواصل الشخصي (الهاتف والرسائل والمراسلات الفورية) في الأعلى؛ يمكن ضبط الأخبار والتسوق ومنصات التواصل الاجتماعي كـ"إشعارات صامتة".
  \item \textbf{تجميع الإشعارات}: اطوِ الإشعارات المتشابهة معاً، لتجنب الغرق في عشرات الرسائل عند فتح القفل مرة واحدة.
\end{itemize}

\booksubsection{الإيماءات السريعة: اجعل عملياتك أسرع من الآخرين}

\begin{itemize}
  \item \textbf{لقطة الشاشة}: زر الطاقة + زر خفض الصوت (iOS بشاشة كاملة)/زر الطاقة + زر خفض الصوت (Android عام).
  \item \textbf{لقطة شاشة طويلة/لقطة شاشة متدفقة}: بعد لقطة الشاشة العادية، اضغط على خيار"لقطة الشاشة الطويلة".
  \item \textbf{مسح الرمز البرمجي بنقرة}: أضف أيقونة"المسح" إلى مركز التحكم/الإعدادات السريعة.
  \item \textbf{النقر على الظهر (iPhone)}: في [الإعدادات] > [إمكانية الوصول] > [اللمس] > [النقر على الظهر]، اربط إجراءً مثل"لقطة الشاشة" أو "قفل الشاشة".
  \item \textbf{التنبيه الصوتي}: Siri/Xiao Yi/Xiaoai/Xiao Bu، نادٍ للاستعلام عن الطقس أو ضبط المنبه.
\end{itemize}

\booksubsection{لوحة المفاتيح: الأداة الأكثر كفاءة والأقل تقديراً}

\begin{itemize}
  \item \textbf{اختيار لوحة المفاتيح المناسبة}:\newline · Sogou/Baidu: قاعدة مفردات كبيرة، تحديث سريع للمصطلحات الشائعة؛\newline · iFlytek: أعلى معدل تعرف على الإدخال الصوتي، أداة رائعة لتدوين الملاحظات في الاجتماعات؛\newline · Gboard (Google): بسيطة، مزامنة قوية عبر الأنظمة، تدعم الإدخال بالانزلاق؛\newline · لوحة iOS المدمجة: أفضل خصوصية، تكامل عميق مع النظام.
  \item \textbf{المفاتيح الضرورية}: التصحيح الذكي، الكلمات التنبؤية، مزامنة قاعدة الكلمات السحابية، الإدخال العلوي للأرقام والرموز.
  \item \textbf{تنبيه خصوصية}: احذر من استخدام وظائف"التنبؤ بالجمل" أو "التوصيات الشخصية"، فقد ترفع محتوى الإدخال.
\end{itemize}

\booksubsection{سطح المكتب وتعدد المهام: اجعل سير عملك سلساً}

\begin{itemize}
  \item \textbf{تقسيم التطبيقات/صورة داخل صورة}: الرد على الرسائل أثناء مشاهدة الفيديو، تدوين الملاحظات أثناء تصفح الويب.
  \item \textbf{النوافذ الصغيرة/النوافذ العائمة}: يدعم WeChat/Douyin"الوضع المصغّر"، أغلق النافذة بعد الرد.
  \item \textbf{تثبيت التطبيقات الشائعة}: ضع فقط 4 تطبيقات الأكثر تكراراً في شريط Dock (أسفل الشاشة)، لتقليل وقت البحث.
\end{itemize}

\begin{tcbinfo}{قائمة فحص"اليوم الأول من العمل"}
بعد الحصول على هاتف جديد أو تحديث كبير للإصدار، يُستحسن قضاء 10 دقائق لمراجعة هذه الإعدادات:\newline (1) كلمة مرور قفل الشاشة + التعرف البيومتري (2) تجميع الإشعارات + إغلاق إشعارات التسويق (3) حجم الخط/معدل التحديث/الوضع الداكن\newline (4) لوحة المفاتيح/المتصفح/البريد الافتراضي (5) جهات الاتصال للطوارئ/البطاقة الطبية (6) إعدادات النسخ الاحتياطي

\end{tcbinfo}

\booksection{خامساً: إمكانية الوصول والاستيعاب: قوة التكنولوجيا من أجل الخير}

لا تدع اسم“إمكانية الوصول”يخيفك، فهذه الوظائف يمكن للجميع استخدامها. صُممت في البداية للأشخاص ذوي الإعاقة، لكنها مفيدة أيضاً للعاديين — إصابة مؤقتة في اليد وحمل طفل، أو تعطل سماعات الأذن في المترو وعدم القدرة على السماع، أو ارتداء نظارات قراءة لرؤية الخط الصغير... كلها يمكن إنقاذها. في النهاية، يعكس تصميم إمكانية الوصول فكرة\textbf{التكنولوجيا من أجل الخير}: تمكين كل شخص من الاستمتاع بالحياة الرقمية على قدم المساواة.

\booksubsection{“عيون”المستخدمين ضعاف البصر — دليل شامل لقارئ الشاشة}

قارئ الشاشة هو الأداة الأساسية للمستخدمين ضعاف البصر لاستخدام الهاتف. سيقوم بقراءة\textbf{جميع المحتويات المرئية}على الشاشة، ويتحكم المستخدم في الهاتف عبر الإيماءات والصوت. فيما يلي طرق استخدام iOS VoiceOver وAndroid TalkBack.

\begin{tcbinfo}{iOS VoiceOver}
\textbf{مسار التفعيل}: [الإعدادات] > [إمكانية الوصول] > [VoiceOver]، افتح المفتاح. يمكن أيضاً النقر ثلاث مرات على الزر الجانبي (أو زر الشاشة الرئيسية) للتبديل السريع.

\textbf{الإيماءات الأساسية}:

\begin{enumerate}
  \item السحب بإصبع واحد: تصفّح محتوى الشاشة عنصراً عنصراً؛
  \item السحب بإصبعين: التمرير لأعلى/لأسفل/يسار/يمين صفحة كاملة؛
  \item النقر المزدوج بإصبع واحد: تفعيل العنصر الذي عليه التركيز (ما يعادل"النقر")؛
  \item السحب بثلاث أصابع: تبديل الصفحات؛
  \item النقر ثلاث مرات بإصبع واحد: التكبير أو التحديد (مختلف عن النقر الثلاثي العادي).
\end{enumerate}

\textbf{الإعدادات الشخصية}: [VoiceOver] > [الصوت]، يمكن ضبط\textbf{سرعة القراءة}(الافتراضي حوالي 50\%، يمكن رفعها إلى 70-80\% لسلاسة أفضل)، و\textbf{طبقة الصوت}، والتبديل بين حزم صوتية مختلفة. تضم iOS ميزات“الصوت المعزّز” و“وحدة التحكم في الصوت”، الأولى تقلل من أصوات التنفس والوقفات، لتجربة استماع أكثر طبيعية.

\end{tcbinfo}

\begin{tcbinfo}{Android TalkBack}
\textbf{مسار التفعيل}: [الإعدادات] > [إمكانية الوصول] > [TalkBack]، افتح المفتاح. أو اضغط مطولاً على زري رفع وخفض الصوت لمدة 3 ثوانٍ للتفعيل السريع (تدعمه بعض العلامات).

\textbf{الإيماءات الأساسية}:

\begin{enumerate}
  \item السحب يساراً/يميناً بإصبع واحد: تبديل عنصر التركيز؛
  \item النقر المزدوج بإصبع واحد: تفعيل العنصر؛
  \item السحب لأعلى/لأسفل/يسار/يمين بإصبعين: التمرير؛
  \item النقر مرة واحدة بثلاث أصابع: الدخول إلى وضع القراءة (قراءة متواصلة)؛
  \item النقر مرة واحدة بأربع أصابع: إيقاف/استئناف القراءة.
\end{enumerate}

\textbf{الإعدادات الشخصية}: [TalkBack] > [إعدادات تحويل النص إلى كلام]، يمكن ضبط\textbf{سرعة الكلام}و\textbf{طبقة الصوت}، والتبديل بين\textbf{محرك الصوت}(يُستحسن تثبيت محرك Google لتحويل النص إلى كلام، لقراءة أكثر طبيعية).

\end{tcbinfo}

\booksubsection{تكبير الخط الصغير وتوضيح الضبابي — تعديلات بصرية}

\booksubsection{إيماءة التكبير}

\begin{itemize}
  \item \textbf{iOS}: [الإعدادات] > [إمكانية الوصول] > [التكبير]، بعد التفعيل\textbf{انقر على الشاشة بثلاث أصابع}للتكبير. التكبير الافتراضي لكامل الشاشة، مع خيار تكبير النافذة (نافذة تكبير موضعية).\textbf{يمكن ضبط نسبة التكبير بفتح وإغلاق إصبعين (1.5x إلى 15x)}، والنقر المزدوج بثلاث أصابع لتثبيت نسبة التكبير الحالية.
  \item \textbf{أندرويد}: [الإعدادات] > [إمكانية الوصول] > [التكبير]، النقر بثلاث أصابع للتكبير، والنقر مرة أخرى بثلاث أصابع للخروج؛ أو\textbf{النقر ثلاث مرات على الشاشة بثلاث أصابع}للتفعيل المؤقت للتكبير (حتى 8 أضعاف).
\end{itemize}

\booksubsection{نص عالي التباين/نص عريض}

\begin{itemize}
  \item \textbf{iOS}: [الإعدادات] > [إمكانية الوصول] > [حجم النص وعرضه]، فعّل مفاتيح"النص العريض" و"تعزيز التباين" و"تقليل الشفافية" وغيرها. يمكن لـ"العكس الذكي"إبقاء النص أسود على خلفية فاتحة، لتجنب نص أبيض على خلفية بيضاء.
  \item \textbf{أندرويد}: [الإعدادات] > [العرض] > [حجم الخط]، يمكن ضبط أربع درجات"صغير، عادي، كبير، ضخم". تقدم Huawei/Xiaomi إضافياً مفتاحي"عريض" و"نص عالي التباين".
\end{itemize}

\booksubsection{مرشحات الألوان (مناسبة لعمى الألوان)}

\begin{itemize}
  \item \textbf{iOS}: [الإعدادات] > [إمكانية الوصول] > [حجم النص وعرضه] > [مرشحات الألوان]، توفر\textbf{أربعة أوضاع لعمى الألوان الأحمر/الأخضر/الأزرق/الأصفر}، بالإضافة إلى خياري"تدرج الرمادي" و"عكس الألوان". يُستحسن إجراء اختبار إدراك الألوان أولاً (مدمج في النظام) قبل اختيار المرشح المناسب.
  \item \textbf{أندرويد}: [الإعدادات] > [إمكانية الوصول] > [تصحيح الألوان]، يحتوي على\textbf{ثلاثة أوضاع لعمى الألوان}+ مفتاح"عكس الألوان".
\end{itemize}

\booksubsection{وصف الصور بالذكاء الاصطناعي: اجعل الذكاء الاصطناعي“عينيك”}

عندما يصادف قارئ الشاشة صورة، فإن الممارسة التقليدية هي تخطيها أو قراءة كلمة“صورة” فقط. لكن الذكاء الاصطناعي يمكنه الآن\textbf{وصف محتوى الصورة بجملة واحدة}، مما يساعد كثيراً في التعرف على رموز التحقق وقراءة الصور وتحديد القوائم.

\begin{tcbinfo}{دعم وصف الصور بالذكاء الاصطناعي في كل نظام}
\begin{itemize}
  \item \textbf{iOS 15+}: فعّل"وصف الصور" في [الإعدادات] > [إمكانية الوصول] > [VoiceOver] > [وصف صوتي]. سيستدعي النظام\textbf{نموذج ذكاء اصطناعي محلي}لإنشاء وصف نصي للصورة تلقائياً (يتطلب iPhone 12 فأحدث).
  \item \textbf{تطبيق TalkBack الجديد}: يدمج Google TalkBack في الإصدارات الأحدث وظيفة\textbf{الأسئلة والأجوبة حول الصور}، يمكنك أن تسأل"ماذا يوجد في هذه الصورة"، ويجيب الذكاء الاصطناعي فوراً (لا يعتمد بشكل صارم على إصدار نظام Android، يمكن تحديث تطبيق TalkBack لتجربته).
  \item \textbf{تطبيقات الطرف الثالث}: "Seeing AI" (iOS) و"Lookout" (Android) و"Xiao Yi Bang Bang" (Huawei) جميعها مساعدون بصريون ممتازون بالذكاء الاصطناعي، يمكنهم التعرف على الأشياء وقراءة المستندات والتعرف على تعابير الوجه.
\end{itemize}

\end{tcbinfo}

\begin{tcbpractice}{سيناريوهات الاستخدام}
جرّب أن تسأل عن صورة شارع“كم شخصاً في هذه الصورة؟ ماذا يفعلون؟”، أو افتح تطبيق توصيل الطعام ودع الذكاء الاصطناعي يقرأ أسماء الأطباق في صورة القائمة —\textbf{هذا هو“التكنولوجيا التي تجعل الحياة أكثر مساواة”حقاً}.

\end{tcbpractice}

\booksubsection{“آذان”المستخدمين ضعاف السمع — دليل شامل للترجمة الفورية}

الترجمة الفورية (Live Captions) هي وظيفة\textbf{تحويل الصوت إلى نص على مستوى النظام}، مفيدة جداً للمستخدمين ضعاف السمع في البيئات الصاخبة والمكتبات والسيناريوهات الأجنبية.

\begin{tcbinfo}{تفعيل وتحسين الترجمة الفورية}
\textbf{iOS}: [الإعدادات] > [إمكانية الوصول] > [الترجمة الفورية]، افتح المفتاح. سيكتشف النظام جميع أصوات الجهاز (المكالمات والفيديو والبودكاست وFaceTime والملاحظات الصوتية)، ويعرض الترجمة تلقائياً أعلى الشاشة.

\textbf{Android 12+}: اضغط\textbf{زر الصوت}انقر على أيقونة“الترجمة الفورية” في النافذة المنبثقة؛ أو فعّلها في [الإعدادات] > [إمكانية الوصول] > [الترجمة الفورية].

\textbf{نصائح تحسين الدقة}:

\begin{enumerate}
  \item استخدمها قدر الإمكان في\textbf{بيئة هادئة}، بعيداً عن المكيفات وأصوات الآلات؛
  \item أعلى دقة في التعرف عندما يكون المتحدث\textbf{واضح النطق، بسرعة معتدلة}؛
  \item الحوار متعدد الأشخاص يقلل الدقة، يمكن استخدام وضع"عرض شخص واحد فقط"؛
  \item يدعم لغات رئيسية مثل الماندرين الصينية والكانتونية والإنجليزية،\textbf{دقة التعرف على اللغات الأقل انتشاراً منخفضة}؛
  \item وجود موسيقى خلفية أو أصوات بشرية متداخلة أثناء الحديث يجعل التعرف متقطعاً.
\end{enumerate}

\end{tcbinfo}

\booksubsection{قناة مزدوجة في سيناريو المكالمات: تحويل الصوت إلى نص + تحويل النص إلى صوت}

عندما يكون المستخدم\textbf{غير قادر على السماع والكلام في الوقت نفسه}(مثل فقدان الصوت المفاجئ، فترة التعافي من جراحة الحنجرة، التأتأة الشديدة)، يمكن للهاتف أن يعمل“مترجماً”.

\begin{tcbinfo}{الترجمة الفورية للمكالمات (TTY/RTT)}
\textbf{iOS}: [الإعدادات] > [إمكانية الوصول] > [RTT/TTY]، بعد التفعيل ستظهر لوحة مفاتيح على الشاشة أثناء المكالمات، ويتحول صوت الطرف الآخر إلى نص في الوقت الفعلي،\textbf{كما ستُقرأ كلماتك المكتوبة بصوت عالٍ}للطرف الآخر.

\textbf{أندرويد}: [الإعدادات] > [إمكانية الوصول] > [النسخ الفوري] + [إشعارات الصوت]. يحوّل“النسخ الفوري”كلام الطرف الآخر إلى نص في الوقت الفعلي؛ وتحوّل“إشعارات الصوت”إدخالك إلى صوت يُيُشغَّل عندما يتعذر استخدام الصوت.

\textbf{التوافق}:

\begin{itemize}
  \item \textbf{RTT (Real-Time Text)}: يتطلب دعم الطرفين، يُعرض النص أثناء الكتابة دون تأخير، تدعمه معظم المشغلات الرئيسية؛
  \item \textbf{TTY (الطباعة عن بُعد)}: يتطلب جهاز TTY مخصصاً، يتم استبداله تدريجياً بـ RTT؛
  \item وظيفة\textbf{تحويل الصوت إلى نص}في WeChat/QQ بديل أيضاً: اضغط مطولاً على الرسالة الصوتية "تحويل إلى نص".
\end{itemize}

\end{tcbinfo}

\booksubsection{وظائف“صغيرة لكن عملية”للجميع}

\begin{itemize}
  \item \textbf{النقر على الظهر}(iOS 14+): [الإعدادات] > [إمكانية الوصول] > [اللمس] > [النقر على الظهر]، يمكن ربط\textbf{نقر مزدوج/ثلاثي}بإطلاق لقطة الشاشة أو قفل الشاشة أو مركز التحكم وغيرها. مفيد جداً عند حمل الأشياء بكلتا اليدين.
  \item \textbf{اللمس المساعد/الكرة العائمة}: iOS [إمكانية الوصول] > [اللمس] > [اللمس المساعد]، Android [إمكانية الوصول] > [التنقل العائم]. زر افتراضي يحل محل الإيماءات أو الأزرار الفيزيائية، ينقذ الموقف عند تعطل الأزرار أو الاستخدام بيد واحدة.
  \item \textbf{التحكم بالتبديل}: التحكم بالهاتف عبر مفاتيح خارجية أو نقرات مفردة أو حركات الرأس،\textbf{مصمم للمستخدمين ذوي الإعاقات الحركية}. iOS [إمكانية الوصول] > [التحكم بالتبديل]، Android [إمكانية الوصول] > [التحكم بالمفاتيح].
  \item \textbf{التعرف على الأصوات}(iOS): [إمكانية الوصول] > [التعرف على الأصوات]، بعد التفعيل"يستمع"الهاتف إلى أصوات مثل جرس الباب وبكاء الرضيع وإنذار الحريق وصفارات الشرطة،\textbf{وعند التعرف يهتز ويعرض إشعاراً على الشاشة}، قيّم جداً للمستخدمين ضعاف السمع وأثناء النوم.
  \item \textbf{تنبيه وميض LED}: iOS [إمكانية الوصول] > [الصوت/البصري] > [وميض LED للتنبيه]، Android [الإعدادات] > [الإشعارات] > [تنبيه الفلاش]. يومض الفلاش عند المكالمات والإشعارات،\textbf{للإدراك حتى في البيئات الصاخبة دون صوت}.
  \item \textbf{الصوت أحادي القناة}: [الإعدادات] > [إمكانية الوصول] > [الصوت] > [الصوت أحادي القناة]، يدمج القناتين اليمنى واليسرى. لن يفقد المستخدمون المصابون بضعف سمع في أذن واحدة نصف المعلومات الصوتية.
\end{itemize}

\begin{tcbinfo}{فلسفة Apple في تصميم إمكانية الوصول}
تستثمر Apple كثيراً في مجال إمكانية الوصول، وتحدّث وتضيف وظائف مساعدة كل عام. يوجد لدى موقع Apple صفحة مخصصة لإمكانية الوصول، تشرح بالتفصيل كيفية استخدام كل وظيفة. عادة ما تكون وظائف إمكانية الوصول في iOS أكثر اكتمالاً وتجربة أكثر سلاسة من Android، وهذا سبب مهم لاختيار العديد من المستخدمين ذوي الإعاقات لـ Apple.

\end{tcbinfo}

\begin{tcbpractice}{اقضِ 30 ثانية لتجربتها}
اختر وقتاً فراغاً، وقم بالتجارب التالية:

\begin{enumerate}
  \item فعّل\textbf{VoiceOver/TalkBack}، وحاول تصفّح WeChat بالسحب بثلاث أصابع دون الاعتماد على عينيك؛
  \item افتح\textbf{الترجمة الفورية}، وتحدث إلى الهاتف لرؤية دقة التعرف؛
  \item استخدم\textbf{النقر بثلاث أصابع}لتكبير نص إلى 200\%، وتجربة يوميات المستخدمين ضعاف البصر؛
  \item أغلق جميع الوظائف السابقة، وفكّر"لو احتجت هذه الوظائف دائماً، كيف أتمنى أن تُصمَّم".
\end{enumerate}

\end{tcbpractice}
